% Clase 26 — longitud de onda lambda = 2 pi c / omega = c / f para las
% frecuencias de la tabla de clase, frente a tamaños de circuito típicos.
% La aproximación cuasi-estacionaria vale si el circuito es mucho más chico
% que lambda. Escala logarítmica en metros.
\begin{tikzpicture}
  \begin{axis}[
    emfig, width=12cm, height=5.2cm,
    xmode=log, log basis x=10,
    xmin=1e-3, xmax=1e8, ymin=0, ymax=1.4,
    ytick=\empty, grid=major,
    xlabel={longitud (m)},
  ]
    % lambdas
    \addplot[only marks, mark=*, mark size=2.4pt, color=figblue, forget plot]
      coordinates {(5e6,1) (300,1) (3,1) (0.03,1)};
    \node[font=\footnotesize, anchor=south] at (axis cs:5e6,1.08) {red, 60 Hz};
    \node[font=\footnotesize, anchor=south] at (axis cs:300,1.08) {AM};
    \node[font=\footnotesize, anchor=south] at (axis cs:3,1.08) {FM};
    \node[font=\footnotesize, anchor=south] at (axis cs:0.03,1.08) {microondas};
    % tamanos de circuito
    \addplot[only marks, mark=square*, mark size=2.2pt, color=figred, forget plot]
      coordinates {(0.3,0.35) (30,0.35) (1e4,0.35)};
    \node[font=\footnotesize, text=figred, anchor=north] at (axis cs:0.3,0.27) {mesa};
    \node[font=\footnotesize, text=figred, anchor=north] at (axis cs:30,0.27) {antena AM};
    \node[font=\footnotesize, text=figred, anchor=north] at (axis cs:1e4,0.27) {ciudad};
  \end{axis}
\end{tikzpicture}
